\SourceLeaf{524}form is $=\frac12(p-1)p^{\pi-3}$; for each of them, however, $M=2p$. The sum of the terms in $\Pi$ arising from these becomes
\[
\begin{aligned}
&=2p\Sigma r^{\mu pp+kp^3}=2p\Sigma r^{\mu pp}\pdot\Sigma r^{kp^3}\\
&=2p\Sigma r^{\mu pp}\pdot\frac{r^{p^\pi}-1}{r^{p^3}-1}=0
\end{aligned}
\]
provided that $\pi>3$. Similarly, the third, fourth, etc. class is represented by $\mu p^4+kp^5$, $\mu p^6+kp^7$, etc., where for $k$ we must take all integers from $0$ to $p^{\pi-5}-1$, $p^{\pi-7}-1$, etc.; for these, $M=2pp$, $M=2p^3$, etc. The sum of the terms in $\Pi$ arising from the third, fourth, etc. class vanishes provided that $\pi>5$, $\pi>7$, etc.. respectively.
\par We conclude that when $\pi$ is even, in $\Pi$ only the terms corresponding to the residue $0$ remain, and these are $=1$; for these, however, $M=p^{\frac12\pi}$, so that the sum of all terms in $\Pi$ is $=p^{\frac12\pi}$.
\BookRule{.27\textwidth}
